\chapter{Конструкторская часть}

\section{Описание алгоритмов}

Схемы рассматриваемых алгоритмов представлены на рисунках 2.1 и 2.2.

\begin{figure}[!htbp]
	\centering
	\includesvg[width=0.4\textwidth]{lab6_bf}
	\caption{Cхема алгоритма полного перебора.}
\end{figure}
\FloatBarrier

\begin{figure}[!htbp]
	\centering
	\includesvg[width=0.4\textwidth]{lab6_aa}
	\caption{Cхема муравьиного алгоритма.}
\end{figure}
\FloatBarrier

\section{Расчёт трудоёмкости алгоритмов}

С помощью введённой модели вычислений (табл.~\ref{t1}) была определена теоретическая трудоёмкость описанных ранее алгоритмов, результаты представлены в табл.~\ref{t2}.

\begin{table}[!htbp]
	\caption{Цены базовых операций.}
	\label{t1}
	\centering
	\begin{tabularx}{\textwidth}{|*{2}{>{\centering\arraybackslash}X|}}
		\hline
		\textbf{Операции} & \textbf{Цена} \\\hline
		{ +, - , =, ++, $--$ , <<, >> } & 1 \\\hline
		{ *, /, \%, *=, /= } & 2 \\\hline
		{ ==, >=, <=, !=, >, <, \&\&, || } & 1 \\\hline
		{ [], \{\} } & 1\\\hline
	\end{tabularx}
\end{table}

\begin{table}[!htbp]
	\centering
	\caption{Трудоёмкость алгоритмов.}
	\label{t2}
	\small
	\begin{tabularx}{\textwidth}{|p{3cm}|*{2}{X|}}
		\hline
		\textbf{Алгоритм} & \textbf{Трудоёмкость ЛС} & \textbf{Трудоёмкость ХС} \\\hline
		Полный перебор & 
		$ (n - 1)!(2n + 1) + n + 4 $ & 
		$ (n - 1)!(2n + 3) + 2n + 2 $ \\\hline
		Муравьиный & 
		$ (T_{max} - 1)n^{3} + 3T_{max}n + 3T_{max} + 2 $ & 
		$ T_{max}n^{3} + (2T_{max} + 2)n + 2T_{max} + 5 $ \\\hline
	\end{tabularx}
\end{table}
\FloatBarrier

\section*{Вывод}

В данном разделе мной были построены схемы алгоритмов и вычислена их трудоёмкость с помощью введённой модели вычислений.

\clearpage
